\documentclass{article}
\usepackage{tcee-repaso}
\tcetema[Balanza de pagos]{3.B.11}{Balanza de pagos: concepto, medición e interpretación.}
\tcefecha{05/11/2026}
\begin{document}
\cabecerarepaso
\begin{columnas}

% ==========================================================================
\bloque{Esquema del cante}
\esquemacante{%
\apartadocante{Introducción}{2,5}
\apartadocante{I. Concepto y medición}{6,5}
\subcante{1. Concepto: BP, PII y otras variaciones}{2,5}
\subcante{2. Medición: residencia y principios contables}{4}
\apartadocante{II. Estructura e interpretación de las cuentas}{11,5}
\subcante{1. Visión general: la identidad}{1}
\subcante{2. Cuenta corriente}{3,5}
\subcante{3. Cuenta de capital}{1}
\subcante{4. Cuenta financiera, PII y rentas}{3}
\subcante{5. Errores y omisiones}{0,5}
\subcante{6. Otras variaciones de la PII}{1}
\subcante{7. Posición de inversión internacional}{1,5}
\apartadocante{III. Interpretación del saldo agregado}{8,5}
\subcante{Entrada: \autor{Meade}, presentación analítica}{0,5}
\subcante{1. CC $=S-I$}{3}
\subcante{2. Reservas (TARGET) y ajuste}{1,5}
\subcante{3. Enfoque de la hoja de balance}{2}
\subcante{4. Desequilibrios globales}{1,5}
\apartadocante{Conclusión}{1}}
\emph{Ejemplo}: almazara cooperativa de Jaén que vende aceite a un importador de Nueva York. \emph{Pizarra}: panorama de las cuentas a la izquierda todo el cante; identidad de la BP y CC $=S-I$ recuadradas a la derecha; el centro se borra en cada apartado. \emph{Dichos}: «toda moneda tiene dos caras» (partida doble), «de grano en grano hace la gallina el granero» (PII), «las cuentas de la lechera» (déficit por rentas futuras).

% ==========================================================================
\bloque[2,5]{Introducción}
\begin{itemize}
\item \textbf{Enganche}: «\emph{No se puede estirar el brazo más que la manga}». España 2007: déficit corriente 9,4\,\% del PIB, el resto del mundo prestaba la diferencia. Desde 2012 presta: capacidad 4,3\,\% (2025); PIIN de $-97$\,\% (2009) a $-42{,}7$\,\% (2025). Hilo: la almazara de Jaén.
\item \textbf{Relevancia}: \emph{política} (supervisión del FMI, \emph{External Sector Report} \vertema{3.B.22}; PDM de la UE: CC media de 3 años fuera de $[-4,+6]$\,\% o PIIN $<-35$\,\%); \emph{teoría} (base del ajuste e intertemporal \vertema{3.B.12} y de las crisis \vertema{3.B.17}; PIIN deudora, sobre todo en deuda, predice crisis: \autor{Catão}-\autor{Milesi-Ferretti}, 2014); \emph{comparación} (mismo manual del FMI).
\item \textbf{Contextualización}: BP + PII = \concepto{cuentas internacionales}; en España, el Banco de España. Parte exterior de la contabilidad nacional (\autor{Kuznets}, \autor{Stone}, años 30-40) \vertema{3.A.30}. Hoy MBP6 (2009, UE desde datos de 2014); MBP7 (marzo 2025, con el SCN 2025, aplicación a finales de la década). Lectura cambiante: Bretton Woods (reservas) $\to$ libre movilidad (déficit que invierte) $\to$ crisis de los 90 y 2008 (importan los \emph{stocks}).
\item \textbf{Problemática}: ¿qué registran BP y PII y con qué reglas? ¿Qué dice cada cuenta? Si la BP siempre cuadra, ¿qué es un desequilibrio? ¿Es malo un déficit? ¿Cuándo preocupa?
\item \textbf{Estructura}: \hyperref[rep:I]{I}, concepto y medición; \hyperref[rep:II]{II}, estructura y lectura de cada cuenta; \hyperref[rep:III]{III}, interpretación del saldo agregado. De la regla contable al diagnóstico.
\end{itemize}

% ==========================================================================
\bloque[6,5]{I. Concepto y medición\label{rep:I}}
\begin{ideaclave}
Las cuentas internacionales son 3 estados enlazados (BP, PII y otras variaciones) elaborados con los principios de la contabilidad nacional: residencia, partida doble, devengo y precios de mercado.
\end{ideaclave}

\sub{I.1. Concepto\label{rep:I.1}}
\begin{itemize}
\item 3 estados (MBP6, cap.~2): \emph{BP} (flujos: transacciones del periodo); \emph{PII} (saldos en una fecha); \emph{cuenta de otras variaciones} (flujos que no son transacciones).
\item \concepto{Balanza de pagos}: registro sistemático y contable de las transacciones económicas entre residentes y no residentes de un país durante un periodo. \emph{Sistemático}: cuentas normalizadas por el FMI; \emph{contable}: partida doble $\Rightarrow$ \textbf{siempre cuadra en conjunto}, no por cuentas; \emph{transacción}: interacción entre 2 unidades, por mutuo acuerdo o imperativo legal, con intercambio o transferencia de valor; \emph{residencia}, no nacionalidad ni moneda; \emph{periodo}: flujo (año; también trimestre y mes).
\item Equilibrio contable $\neq$ económico: \concepto{desequilibrio exterior} = saldo parcial (CC o capacidad/necesidad) que no se puede mantener indefinidamente.
\item \concepto{PII}: balance financiero con el exterior en una fecha. \concepto{PIIN} $\equiv$ activos exteriores $-$ pasivos exteriores: $>0$ acreedor neto, $<0$ deudor neto. Parte exterior del balance de la nación.
\item \concepto{Otras variaciones}: revalorizaciones (precios, tipo de cambio) y otros cambios de volumen; fuera de la BP, pero sin ellas no se explica la PII final.
\end{itemize}
\graficorepaso[0.97\linewidth]{panorama-cuentas-internacionales}{Panorama de las cuentas internacionales: PII inicial + CF + OV = PII final}
\begin{itemize}
\item Almazara: vende a Nueva York $\Rightarrow$ BP; a Sevilla $\Rightarrow$ no aparece; depósito en dólares en Nueva York $\Rightarrow$ activo de la PII; si el dólar se deprecia, vale menos euros sin transacción $\Rightarrow$ otra variación.
\end{itemize}

\sub{I.2. Medición\label{rep:I.2}}
\begin{itemize}
\item Manual del FMI: 1.\textsuperscript{a} ed. 1948; MBP5 1993; \textbf{MBP6 2009}; MBP7 20/3/2025 (aplicación 2029-2030). Cada revisión: más cerca del SCN y más peso de los \emph{stocks}.
\item \emph{MBP5 $\to$ MBP6}: PII en el título; \textbf{signo de la CF} = VNA $-$ VNP; rentas y transferencias $\to$ rentas primarias y secundarias; transformación sin cambio de propiedad $\to$ servicios; \emph{merchanting} $\to$ bienes; I+D, activo producido (fuera de la CK). \emph{MBP7}: criptoactivos sin pasivo de contrapartida (bitcóin) = activo no financiero no producido en la CK, no reservas; reservas internacionales netas; más detalle sectorial y digital.
\item \concepto{Flujos}: \emph{transacciones} (intercambios o transferencias) y \emph{otros flujos} (revalorización, cancelación unilateral, catástrofe, expropiación sin indemnización). \concepto{Posiciones}: saldos en una fecha, unidos por los flujos.
\item \concepto{Residente}: \emph{centro de interés económico predominante} en el territorio económico (referencia: 1 año o más, no rígida). Hogares: vivienda principal (estudiantes, pacientes y turistas, residentes de origen); filiales donde operan; sucursales si producción significativa $\geq$1 año; entidades con fines especiales, donde se constituyen; AA.~PP. con embajadas y bases; \textbf{organismos internacionales: de ninguna economía} (el BCE es resto del mundo para España). Filial de la almazara en Nueva Jersey: residente de EE.~UU.; el capital, inversión directa.
\item \emph{4 sectores} (MBP6): banco central; sociedades de depósito salvo el BC (otras IFM); AA.~PP.; otros sectores. La CN (SEC 2010): 5 (S.11-S.15), con el BC dentro de S.12.
\item \concepto{Partida doble} (vertical): cada transacción, 2 asientos por igual importe. \emph{Créditos}: X, rentas a cobrar, $\downarrow$ activos, $\uparrow$ pasivos. \emph{Débitos}: M, rentas a pagar, $\uparrow$ activos, $\downarrow$ pasivos. CC y CK en bruto; CF en neto por instrumento: \concepto{VNA} y \concepto{VNP}, $+$ si aumentan. \emph{Horizontal}: las 2 partes, mismo importe y momento. \concepto{Partida cuádruple}: 4 asientos en el sistema mundial (en la práctica no se cumple: el mundo no cuadra).
\item Almazara: exporta 100 y cobra en Nueva York (bienes $+100$; VNA $+100$); vende acciones por 50 a depósito (solo composición); préstamo italiano de 70 para una envasadora (bienes $-70$; VNP $+70$) $\Rightarrow$ CC $=30=$ CF.
\item \concepto{Devengo}: cuando se crea, transforma, intercambia, transfiere o extingue el valor, no al cobrar. Bienes y activos: cambio de \emph{propiedad económica}; servicios: al prestarse; rentas: al nacer el derecho (intereses, día a día). Pago aplazado $\Rightarrow$ \concepto{crédito comercial} en la CF.
\item \emph{Valoración}: precios de mercado (préstamos y depósitos a valor nominal; no cotizadas, por fondos propios); moneda nacional al tipo del momento (posiciones, al del día de referencia). \textbf{Mercancías FOB} en X y M; Aduanas (ONU): M \concepto{CIF} $\Rightarrow$ M de Aduanas $>$ M de la BP.
\item BP frente a CN: mismos principios; distinto \emph{punto de vista}: la cuenta del resto del mundo se ve desde fuera, signo cambiado. Saldo común: CC $+$ CK = B.9.
\item \emph{Fuentes} (BdE): proveedores de servicios de pago ($>$50.000~€; común a CC y CF); Aduanas (DUA e Intrastat, ajustadas: CIF$\to$FOB, \emph{merchanting}, transformación); INE (EGATUR, FRONTUR, servicios); declaraciones directas (Circular 4/2012: mensual $\geq$300~M€); Registro de Inversiones Exteriores. Reglamento (CE) 184/2005, adaptado por el (UE) 555/2012.
\end{itemize}

% ==========================================================================
\bloque[11,5]{II. Estructura e interpretación de las cuentas\label{rep:II}}
\begin{ideaclave}
CC y CK dicen cuánto presta o pide prestado el país; la CF, con qué instrumentos y qué sectores; la PII, cuánto ha acumulado y con qué riesgo.
\end{ideaclave}

\sub{II.1. Visión general de las cuentas\label{rep:II.1}}
\begin{itemize}
\item 3 cuentas: \concepto{cuenta corriente} (CC: bienes y servicios, RP, RS); \concepto{cuenta de capital} (CK: transferencias de capital, activos no financieros no producidos); \concepto{cuenta financiera} (CF $=$ VNA $-$ VNP). Las 2 primeras, créditos y débitos; la CF, VNA y VNP. Todo desde el residente. Con $\mathrm{CC}=(X-M)+\mathrm{RP}+\mathrm{RS}$, \textbf{identidad} (MBP6):
\funcionclave{\mathrm{CC}+\mathrm{CK}+\mathrm{EO}=\mathrm{CF}=\mathrm{ID}+\mathrm{IC}+\mathrm{D}+\mathrm{OI}+\Delta R}
\item CC $+$ CK $>0$: \concepto{capacidad de financiación} (acumula activos o cancela pasivos: sale capital neto, CF $>0$); $<0$: \concepto{necesidad de financiación} (se endeuda o vende activos: entra capital).
\item España 2025: capacidad 4,3\,\% del PIB (servicios compensan con creces bienes; CK, 1 punto por fondos UE); 12 meses a agosto de 2026: 3,7\,\%.
\end{itemize}

\sub{II.2. Cuenta corriente\label{rep:II.2}}
\begin{itemize}
\item \emph{Bienes}: cambio de propiedad económica, crucen o no la frontera: mercancías generales, \emph{merchanting} (por el margen), oro no monetario. \emph{Servicios}: 12 componentes (transporte, viajes, financieros, propiedad intelectual, otros empresariales\ldots). Almazara: aceite = bienes; flete de naviera española = servicios; canon de marca extranjera = M de servicios.
\item Servicios: algo más de 1/4 del comercio mundial bruto (9,6 de 34,7 billones de \$, 2025); en valor añadido $\approx$1/2: \concepto{servicificación} (OCDE, TiVA). \emph{Curva de la sonrisa} (\autor{Shih}): más valor añadido en I+D, diseño, marca y posventa \vertema{3.B.18}.
\item Saldo: $-$ renta nacional, $+$ renta de los socios, $+$ competitividad. \emph{Precio}: \concepto{tipo de cambio efectivo real},
\begin{formula}
\centering $q^{ef}=\prod_i\left(\dfrac{E_iP_i^*}{P}\right)^{\omega_i},\quad\sum_i\omega_i=1$
\end{formula}
\item[] $\uparrow q^{ef}$ = depreciación real = ganancia de competitividad \vertema{3.B.13}; deflactores: IPC, CLU o precios de exportación. \textbf{Ojo}: BCE y BdE publican al revés (subida = apreciación). \emph{No precio}: capital físico, humano, tecnológico, calidad, marca, instituciones (B-READY sustituye al \emph{Doing Business}, 2021). \emph{Resultados}: cuota de mercado, tasa de cobertura $X/M$, propensión exportadora $X/\mathrm{PIB}$, contenido importado.
\item \emph{Unión monetaria}: $\gamma_q\simeq\gamma_E+\pi^*-\pi$ y, con $E$ fijo,
\begin{formula}
\centering $\gamma_q\simeq\pi^*-\pi$: inflación $>$ socios = apreciación real \vertema{3.A.41}
\end{formula}
\item[] CLU: $\gamma_{\mathrm{CLU}}=\gamma_W-\gamma_{\mathrm{PAT}}$. Solo queda la \concepto{devaluación interna} (lenta, costosa en empleo; sur del euro tras la crisis \vertema{3.B.16}, \vertema{4.A.22}). Pero \autor{Balassa}-\autor{Samuelson} (1964): país que converge $\Rightarrow$ inflación de no comerciables sin pérdida de competitividad $\Rightarrow$ mirar CLU de manufacturas o precios de exportación. El PDM vigila TCER y CLU.
\item España 2025: bienes $-3{,}0$\,\% (energía); servicios $+6{,}8$\,\%, turismo 4,3 puntos.
\item \concepto{Rentas primarias}: uso temporal de factores: \emph{trabajo} (fronterizos, temporada, $<$1 año), \emph{inversión} (intereses, dividendos, beneficios reinvertidos; el grueso), \emph{otra} (impuestos sobre producción e importaciones, subvenciones, recursos naturales). Explican PIB $\neq$ RNB (Irlanda: RNB $\ll$ PIB). Más deudor $\Rightarrow$ más rentas pagadas $\Rightarrow$ peor CC futura.
\item \concepto{Rentas secundarias}: transferencias corrientes: remesas, cooperación corriente, aportaciones UE (IVA, RNB), impuestos sobre la renta y prestaciones, seguros no vida. \emph{NGEU}: gasto corriente $\to$ RS; inversión $\to$ CK; préstamos del MRR $\to$ OI de AA.~PP. (intereses, RP).
\end{itemize}

\sub{II.3. Cuenta de capital\label{rep:II.3}}
\begin{itemize}
\item \emph{Transferencias de capital} (activo fijo, ayuda condicionada a invertir, condonación pactada; España: FEDER, Cohesión, MRR de inversión) y \emph{compraventa de activos no financieros no producidos} (tierra, pesca, espectro, contratos, licencias, marcas, logotipos, dominios).
\item \textbf{No va aquí}: patentes y derechos de autor (producidos desde el SCN 2008 $\to$ servicios); cancelación unilateral (otro cambio de volumen); transferencias de gasto corriente (RS).
\item CC $+$ CK = capacidad/necesidad; CK pequeña en avanzados, no en España desde 2021 (1,0\,\% en 2025). Almazara: ayuda UE para la envasadora = CK $+$; la máquina = bienes $-$.
\item \textbf{Cuenta de capital $\neq$ «balanza de capital»} de los manuales antiguos (= hoy CF).
\end{itemize}

\sub{II.4. Cuenta financiera, PII y rentas de la inversión\label{rep:II.4}}
\begin{itemize}
\item 3 caras: CF (flujo), PII (saldo), rentas de la inversión (rendimiento). Encadenadas: $\uparrow$ entradas netas $\Rightarrow$ $\downarrow$ PIIN $\Rightarrow$ $\uparrow$ rentas pagadas $\Rightarrow$ $\downarrow$ CC. Se netea por instrumento, no activos con pasivos; clasificación \emph{funcional}:
\begin{enumerate}
\item \concepto{Inversión directa}: control o influencia significativa, $\geq$10\,\% de los votos; participaciones, beneficios reinvertidos, deuda intragrupo. La más estable; trae tecnología \vertema{4.A.23}.
\item \concepto{Inversión de cartera}: acciones $<$10\,\%, fondos, valores de deuda negociables; sin control, volátil.
\item \emph{Derivados} (y opciones de empleados): solo neto \vertema{3.B.25}.
\item \emph{Otra inversión} (residual): préstamos, depósitos, créditos comerciales, repos, billetes; en el euro, posición del BCN frente al Eurosistema.
\item \concepto{Reservas internacionales}: activos exteriores que el BC controla y puede usar de inmediato: oro monetario, DEG, posición en el FMI, divisas. $\Delta R>0$ si aumentan.
\end{enumerate}
\item Compra de divisas $\Rightarrow$ $\uparrow$ base monetaria; \concepto{esterilizar}: vender títulos en moneda nacional \vertema{3.B.15}. Motivos: \emph{precaución} (parada súbita; tras 1997), \emph{frenar la apreciación} (China, 2000), materias primas (fondos soberanos). \concepto{Dilema de \autor{Triffin}} (1960): el mundo acumula dólares si EE.~UU. tiene déficit, que erosiona la confianza.
\item Reservas de divisas: 13,2 billones de \$ (jun. 2026), China 1/4; dólar 71\,\% (1999) $\to$ 56,7\,\%; euro 20,6\,\%; yen 5\,\%; libra 4,6\,\%.
\item \emph{Zona del euro}: reservas, activos netos frente al Eurosistema, otros. La posición frente al Eurosistema cuadra la BP:
\begin{formula}
\centering $\mathrm{CF}=\mathrm{CF}^{\text{exc. BdE}}+\mathrm{CF}^{\text{BdE}},\ \mathrm{CF}^{\text{BdE}}=\Delta R+\Delta\text{TARGET}+\text{otros}$
\end{formula}
\item[] T2 sustituye a TARGET2 (20/3/2023). Importador francés paga a la almazara: $\uparrow$ activo del BdE frente al Eurosistema; saldos suman 0. Saldo deudor que crece: parada súbita (\autor{Sinn}-\autor{Wollmershäuser}, 2012: rescate encubierto) o reparto de la liquidez y del programa de compras (\autor{Whelan}, 2014). BdE: $\approx-414.000$~M€ (fin 2025) \vertema{4.A.21}.
\end{itemize}

\sub{II.5. Errores y omisiones\label{rep:II.5}}
\begin{itemize}
\item Los 2 asientos salen de fuentes distintas (Aduanas, bancos) $\Rightarrow$ residuo:
\begin{formula}
\centering $\mathrm{EO}=\mathrm{CF}-(\mathrm{CC}+\mathrm{CK})$
\end{formula}
\item[] EO $>0$: CF sobrevalorada o CC $+$ CK infravaloradas. Grandes, persistentes y de un signo: fuga de capital no declarada. \textbf{No son una cuenta}; no se reparten. España 2025: 0,8\,\% del PIB.
\end{itemize}

\sub{II.6. Cambios en la PII no recogidos en la CF\label{rep:II.6}}
\begin{itemize}
\item Variación de la posición:
\begin{formula}
\centering $\mathrm{PIIN}_t-\mathrm{PIIN}_{t-1}=\mathrm{CF}_t+\mathrm{OV}_t$
\end{formula}
\item[] OV: \emph{revalorizaciones} (precios, tipo de cambio) y \emph{otros cambios de volumen} (cancelaciones unilaterales, reclasificaciones, expropiaciones, cambios de residencia).
\item \emph{Depreciación}: $\uparrow$ valor de lo que está en divisas: mejora a EE.~UU. (pasivos en \$), empeora a emergentes con deuda en \$. Bonos vuelven al nominal; el efecto se concentra en acciones. \textbf{Lo que pesa}: deuda en divisas revalorizada por una depreciación (\autor{Lane}-\autor{Milesi-Ferretti}, 2007).
\item España 2025: PIIN $-83.000$~M€ pese a CF $+85.200$: OV $\approx-168.000$ ($\approx$10 puntos del PIB).
\end{itemize}

\sub{II.7. Posición de inversión internacional\label{rep:II.7}}
\begin{itemize}
\item Balance frente al exterior, mismas categorías funcionales que la CF; suma de CF pasadas $+$ OV: «\emph{de grano en grano hace la gallina el granero}».
\item 3 medidas: \concepto{PIIN}; \concepto{PIIN excluidos los instrumentos sin riesgo de impago} (Eurostat: sin acciones ni deuda intragrupo); \concepto{deuda externa bruta} (solo pasivos con pago de principal o intereses).
\item Lectura: signo y nivel (PDM: $-35$\,\%); por qué cambia; \textbf{composición}: capital mejor que deuda (y ID), largo mejor que corto, moneda nacional, sectores (pasivo del BC frente al Eurosistema $\neq$ deuda bancaria a corto). ¿Brutos (\autor{Shin}, 2012) o netos de deuda (\autor{Catão}-\autor{Milesi-Ferretti}, 2014)?
\item España: $-97{,}0$\,\% (2009) $\to$ $-39{,}9$\,\% (2024) $\to$ $-42{,}7$\,\% (2025, revalorizaciones) $\to$ $-40{,}7$\,\% (mediados de 2026); sin BdE $-31{,}1$\,\%; sin instrumentos sin riesgo de impago, $-9{,}6$\,\% (2025); deuda externa bruta 168\,\%.
\end{itemize}

% ==========================================================================
\bloque[8,5]{III. Interpretación del saldo agregado\label{rep:III}}
\begin{ideaclave}
Un saldo exterior no es bueno ni malo en sí: el saldo corriente es la brecha ahorro-inversión; se financia con capitales o reservas (TARGET en el euro) y su sostenibilidad depende de la posición que deja, su composición y el crecimiento.
\end{ideaclave}
\begin{itemize}
\item \autor{Meade} (1951): transacciones \emph{autónomas} (por sus motivos) frente a \emph{compensatorias} (de las autoridades: reservas, FMI); desequilibrio = hacen falta las segundas. MBP6: \concepto{presentación analítica}: por encima de la línea todo salvo reservas y vinculadas = \emph{saldo global}; debajo, $\Delta R$, crédito del FMI y financiación excepcional.
\end{itemize}

\sub{III.1. Cuenta corriente: brecha entre ahorro e inversión\label{rep:III.1}}
\begin{itemize}
\item $\mathrm{RNDB}=\mathrm{PIB}+\mathrm{RP}+\mathrm{RS}$; absorción $A=C+I+G$; $Y=A+(X-M)$; $S=\mathrm{RNDB}-C-G$:
\funcionclave{\mathrm{CC}=(Y-A)+\mathrm{RP}+\mathrm{RS}=S-I=(S_P-I_P)+(T-G)}
\item Déficit = gastar más de lo que se ingresa o invertir más de lo que se ahorra. \textbf{Identidades, no teorías}: un arancel o una devaluación solo mueven la CC si cambian $S$ o $I$ \vertema{3.B.12}.
\item \emph{Mismo déficit, distinto diagnóstico}: poco ahorro (sobrecalentamiento) o inversión productiva (convergencia); construcción con crédito: España 2000-2007. «\emph{Las cuentas de la lechera}». Corto plazo: igual ahorrar que invertir menos; largo, no.
\end{itemize}
\graficorepaso[0.97\linewidth]{arbol-deficit-corriente}{Diagnóstico de un déficit: financiar lo temporal, ajustar lo permanente}
\begin{itemize}
\item \concepto{Déficits gemelos}: déficit público $\Rightarrow$ corriente (multiplicador; \autor{Mundell}-\autor{Fleming}: tipos, entradas, apreciación). Equivalencia ricardiana: no \vertema{3.A.39}. Evidencia: consolidación de 1 punto $\Rightarrow$ CC $+0{,}6$ (\autor{Bluedorn}-\autor{Leigh}, 2011).
\item \concepto{Enfoque intertemporal} (\autor{Sachs}, 1981; \autor{Obstfeld}-\autor{Rogoff}, 1995): CC = reparto de consumo e inversión en el tiempo; pedir prestado ante un choque temporal o para invertir donde rinde más es óptimo si la deuda es sostenible.
\item \concepto{Paradoja de \autor{Feldstein}-\autor{Horioka}} (1980): OCDE, $I$ sigue a $S$ (coeficiente $\approx$0,9), contra la movilidad perfecta; sesgo nacional, costes, riesgo cambiario, objetivo exterior. Cae con la integración, sobre todo en el euro (Grecia, Portugal: \autor{Blanchard}-\autor{Giavazzi}, 2002).
\end{itemize}

\sub{III.2. Reservas y ajuste de la balanza de pagos\label{rep:III.2}}
\begin{itemize}
\item Despejando la identidad:
\begin{formula}
\centering $\Delta R=(\mathrm{CC}+\mathrm{CK})+\mathrm{EO}-(\mathrm{ID}+\mathrm{IC}+\mathrm{D}+\mathrm{OI})$
\end{formula}
\item[] Las reservas cubren lo que no cubre el capital privado; pueden caer con capacidad si sale más capital. Tipo fijo: reservas (salvo esterilización, $\downarrow$ dinero); flexible: tipo de cambio \vertema{3.B.15}.
\item Euro: $\Delta\text{TARGET}+\Delta R+\text{otros}=\mathrm{CF}^{\text{BdE}}$; \textbf{TARGET sin límite}: ajuste por el coste de la financiación privada, no por el BC.
\item Financiación cortada: temporal $\Rightarrow$ reservas suavizan; permanente $\Rightarrow$ ataque especulativo \vertema{3.B.17}; sin reservas ni crédito: \concepto{parada súbita} (\autor{Calvo}, 1998).
\item \emph{Ajuste}: reducir el gasto (fiscal y monetaria) y desviarlo (depreciación si \autor{Marshall}-\autor{Lerner}: $|\varepsilon_X|+|\varepsilon_M|>1$), reformas, financiación oficial; sin ajuste del gasto, la depreciación se va en inflación.
\item \emph{Suficiencia de reservas} (FMI, 2011): 3 meses de M; 100\,\% de la deuda a corto (\autor{Greenspan}-\autor{Guidotti}); 20\,\% de M2; métrica ARA. Demasiadas también cuestan.
\end{itemize}

\sub{III.3. El enfoque de la hoja de balance\label{rep:III.3}}
\begin{itemize}
\item \textbf{a) Idea}: las crisis de los 90 (Asia, Rusia, Brasil, Argentina) no se veían en la CC, sino en los \emph{balances} por sectores (\autor{Allen} et al., 2002).
\item \textbf{b) Modelo}. \emph{Supuestos}: precios de mercado; consolidado, se cancela lo de residentes. \emph{Desarrollo}: riqueza nacional neta $=K+\mathrm{PIIN}$;
\begin{formula}
\centering $\Delta\text{riqueza}=(S-\delta K)+\mathrm{CK}+\mathrm{OV}^K+\mathrm{OV}^F$
\end{formula}
\item[] Dinámica en \% del PIB, con $\gamma^N$ el crecimiento nominal:
\begin{formula}
\centering $\mathrm{PIIN}_t-\mathrm{PIIN}_{t-1}\approx\mathrm{CF}_t+\mathrm{OV}_t-\dfrac{\gamma^N_t}{1+\gamma^N_t}\mathrm{PIIN}_{t-1}$
\end{formula}
\end{itemize}
\graficorepaso[0.97\linewidth]{balance-nacion}{El balance de la nación y sus descalces}
\begin{itemize}
\item \emph{Implicaciones}: solo enriquece ahorrar, recibir transferencias de capital o ganar plusvalías; pedir prestado no (si financia inversión, cambia deuda por capital; si consumo, empobrece). \emph{Sostenibilidad}: estable si el saldo sin rentas de la inversión cubre tipo $-$ crecimiento: la manga es el crecimiento; el brazo, el tipo. Saldo estabilizador (anexo):
\begin{formula}
\centering $b_t=\dfrac{1+i^*}{1+\gamma^N}b_{t-1}+d_t-\mathrm{ov}_t\ \Rightarrow\ -d=\dfrac{i^*-\gamma^N}{1+\gamma^N}\,b$
\end{formula}
\item[] $i^*>\gamma^N$: superávit solo para estabilizar (\emph{bola de nieve}). Con $b=41$\,\% y $\gamma^N-i^*=2$ puntos, cabe un déficit de $\approx$0,8 puntos.
\item \emph{Descalces}: de \emph{estructura de capital} (deuda, no acciones); de \emph{plazos}; de \emph{monedas} (\concepto{pecado original}, \autor{Eichengreen}-\autor{Hausmann}-\autor{Panizza}, 2005; crisis gemelas, \autor{Kaminsky}-\autor{Reinhart}, 1999); de \emph{solvencia} (contagio entre sectores vía banca); \emph{concentración} de acreedores. Equilibrios múltiples (\autor{Obstfeld}, 1996) $\Rightarrow$ margen.
\item \textbf{c) Valoración}: completa los flujos con los saldos; FMI y PDM lo usan. Límites: datos por sector, moneda y plazo; no dice qué deuda es demasiada.
\end{itemize}

\sub{III.4. Desequilibrios globales\label{rep:III.4}}
\begin{itemize}
\item \concepto{Desequilibrios globales}: saldos corrientes grandes y persistentes de las grandes economías; suma de valores absolutos en \% del PIB mundial (los saldos mundiales no suman 0 por medición).
\item Máximo 2006-2008 (EE.~UU. frente a China, Japón, Alemania, petroleros); se amplían: 4,0 $\to$ 4,4 billones de \$ (2024-2025), $\approx$3,7\,\% del PIB mundial; FMI: 2/3 excesivos. 2025: EE.~UU. $-3{,}7$\,\%; China $+3{,}7$, Alemania $+4{,}5$, Japón $+4{,}9$.
\item \emph{Por qué}: \concepto{exceso de ahorro mundial} (\autor{Bernanke}, 2005); \concepto{privilegio exorbitante}: de banquero a capitalista de riesgo del mundo (\autor{Gourinchas}-\autor{Rey}, 2007); causas internas (déficit público de EE.~UU.; ahorro alto e inversión débil en China y Alemania).
\item \emph{¿Problema?} No si son decisiones intertemporales (envejecimiento); sí si nacen de distorsiones o acumulan riesgos (\autor{Blanchard}-\autor{Milesi-Ferretti}, 2012). Causas comunes con 2008 (\autor{Obstfeld}-\autor{Rogoff}, 2009).
\item \concepto{Materia oscura} (\autor{Hausmann}-\autor{Sturzenegger}, 2007): EE.~UU. deudor que cobra más rentas; discutida (\autor{Curcuru}-\autor{Dvorak}-\autor{Warnock}, 2008).
\item Vigilancia: FMI (metodología EBA, ESR) \vertema{3.B.22}; PDM; G-20 \vertema{3.B.19}. España: superávit desde 2012 (3,2\,\% en 2025), dentro de la banda \vertema{4.A.21}.
\end{itemize}

% ==========================================================================
\bloque[1]{Conclusión}
\begin{itemize}
\item \textbf{Recapitulación}: en conclusión, la BP registra las transacciones con el exterior con las reglas de la contabilidad; la PII, lo acumulado; las OV las unen. Siempre cuadra: importan los saldos parciales, sobre todo capacidad/necesidad = $S-I$. Un déficit no es bueno ni malo: origen, financiación y sostenibilidad.
\item \textbf{Relevancia}: base de la supervisión del FMI y del PDM; parte del dictamen de coyuntura.
\item \textbf{Extensiones}: ajuste e intertemporal \vertema{3.B.12}, crisis \vertema{3.B.17}, capitales \vertema{3.B.18}, FMI y reservas \vertema{3.B.22}, BP de España \vertema{4.A.21}; reto del MBP7: digitalización, criptoactivos, propiedad intelectual de multinacionales.
\item \textbf{Opinión}: en una unión monetaria la BP no desaparece: se esconde en TARGET; mirar menos el saldo y más origen, financiación y posición con su composición.
\item \textbf{Idea final}: «no se puede estirar el brazo más que la manga», salvo que otro preste la manga: la BP dice cuánto se estira y quién presta; la PII, hasta cuándo.
\end{itemize}

% ==========================================================================
\bloque{Autores y fechas}
\hito{1948}{1.\textsuperscript{er} Manual de Balanza de Pagos del FMI}
\hito{1951}{\autor{Meade}: transacciones autónomas y compensatorias}
\hito{1960}{\autor{Triffin}: dilema del dólar}
\hito{1964}{\autor{Balassa}-\autor{Samuelson}}
\hito{1980}{\autor{Feldstein}-\autor{Horioka}: paradoja}
\hito{1981}{\autor{Sachs}: enfoque intertemporal}
\hito{1993}{MBP5}
\hito{1995}{\autor{Obstfeld}-\autor{Rogoff}: enfoque intertemporal}
\hito{1998}{\autor{Calvo}: paradas súbitas}
\hito{1999}{\autor{Kaminsky}-\autor{Reinhart}: crisis gemelas}
\hito{2002}{\autor{Allen} et al.: enfoque de balances; \autor{Blanchard}-\autor{Giavazzi}}
\hito{2005}{\autor{Bernanke}: exceso de ahorro; pecado original}
\hito{2007}{\autor{Gourinchas}-\autor{Rey}; \autor{Hausmann}-\autor{Sturzenegger}; \autor{Lane}-\autor{Milesi-Ferretti}}
\hito{2009}{MBP6}
\hito{2012}{\autor{Sinn}-\autor{Wollmershäuser}: TARGET; Reglamento 555/2012}
\hito{2014}{\autor{Catão}-\autor{Milesi-Ferretti}; MBP6 en la UE}
\hito{2023}{T2 sustituye a TARGET2}
\hito{2025}{MBP7 y SCN 2025}

% ==========================================================================
\bloque{Datos clave}
\begin{itemize}
\item España: capacidad de financiación \dato{4,3\,\% del PIB (2025)}{BdE, 2026}; CC \dato{3,2\,\%}{BdE, 2026}; bienes \dato{$-3{,}0$\,\%}{BdE, 2026}; servicios \dato{$+6{,}8$\,\%, turismo 4,3}{BdE, 2026}; EO \dato{0,8\,\%}{BdE, 2026}.
\item PIIN: \dato{$-42{,}7$\,\% del PIB (2025)}{BdE, 2026}; \dato{$-40{,}7$\,\% (mediados de 2026)}{BdE, 2026}; mínimo \dato{$-97$\,\% (2009)}{BdE, 2026}; deuda externa bruta \dato{168\,\%}{BdE, 2026}.
\item Saldo TARGET del BdE: \dato{$\approx-414.000$~M€ (fin 2025)}{BCE, 2026}.
\item Reservas de divisas mundiales: \dato{13,2 billones de \$; dólar 56,7\,\%, euro 20,6\,\%}{FMI COFER, jun. 2026}.
\item Desequilibrios: \dato{4,4 billones de \$, $\approx$3,7\,\% del PIB mundial}{FMI ESR, 2026}.
\end{itemize}

% ==========================================================================
\bloque{Preguntas del tribunal}
\begin{enumerate}
\item \textbf{Si la BP siempre cuadra, ¿qué es un desequilibrio?} Cuadra por partida doble; desequilibrio = saldo parcial insostenible o necesidad de transacciones compensatorias (\autor{Meade}): saldo global de la presentación analítica.
\item \textbf{¿Es malo un déficit corriente?} No en sí: si financia inversión productiva o suaviza un choque temporal (enfoque intertemporal) es óptimo; preocupa si es por poco ahorro, consumo o crédito y deja una PIIN grande en deuda, a corto y en divisas.
\item \textbf{¿Por qué empeoró la PIIN española en 2025 con CF positiva?} Revalorizaciones ($\approx-168.000$~M€): subieron de precio los pasivos (acciones españolas); no son más obligaciones de pago.
\item \textbf{¿Qué es TARGET y qué indica?} Posición del BCN frente al Eurosistema; hace de reservas sin límite; deudor creciente = parada súbita privada (\autor{Sinn}) o reparto de la liquidez (\autor{Whelan}).
\item \textbf{¿Qué medida de la PII predice mejor las crisis?} Netos de deuda (\autor{Catão}-\autor{Milesi-Ferretti}) o brutos (\autor{Shin}); se mira también la PIIN sin instrumentos sin riesgo de impago y la deuda externa bruta.
\item \textbf{¿Qué cambia el MBP7?} Criptoactivos sin pasivo a la CK, no reservas; reservas netas normalizadas; más detalle sectorial y digital; se aplicará a finales de la década.
\item \textbf{¿Causa el déficit público el corriente?} Déficits gemelos (multiplicador; \autor{Mundell}-\autor{Fleming}); con equivalencia ricardiana, no; evidencia: 1 punto de consolidación $\Rightarrow$ CC $+0{,}6$.
\item \textbf{¿Cuántas reservas bastan?} 3 meses de M, 100\,\% de la deuda a corto (\autor{Greenspan}-\autor{Guidotti}), 20\,\% de M2 o la métrica ARA del FMI; tenerlas y esterilizarlas cuesta.
\item \textbf{¿Cuándo es sostenible una PIIN deudora?} Si el saldo sin rentas de la inversión cubre $(i^*-\gamma^N)\,b/(1+\gamma^N)$; importan también composición (deuda, plazo, moneda) y efectos de valoración.
\end{enumerate}

\begin{trampa}
\begin{itemize}
\item MBP6: CF $=$ VNA $-$ VNP; capacidad (CC $+$ CK $>0$) = CF $>0$ (salida neta). MBP5: signo contrario, CC $+$ CK $+$ CF$'$ $+$ EO $=0$.
\item Créditos: X, rentas a cobrar, $\downarrow$ activos, $\uparrow$ pasivos; débitos: M, $\uparrow$ activos, $\downarrow$ pasivos.
\item Cuenta de capital $\neq$ «balanza de capital» antigua (= CF). Patentes: servicios; marcas y dominios: CK.
\item Residencia, no nacionalidad: organismos internacionales, de ninguna economía; estudiantes y pacientes, del país de origen; el BCE es resto del mundo para España.
\item Mercancías FOB en X y M en la BP; Aduanas: M CIF.
\item ID: $\geq$10\,\% de los votos. EO: residuo, no una cuenta.
\item TCER: $\uparrow q^{ef}$ = depreciación real; BCE y BdE publican al revés.
\item Beneficios reinvertidos: renta primaria y, a la vez, ID. NGEU: corriente $\to$ RS; inversión $\to$ CK.
\item Compra de divisas por el BC: $\uparrow$ base monetaria; esterilizar = \emph{vender} títulos en moneda nacional.
\item \autor{Feldstein}-\autor{Horioka}: correlación $S$-$I$ alta ($\approx$0,9) = \emph{poca} movilidad aparente del capital.
\item Depreciación: mejora la PIIN de quien tiene activos en divisas y pasivos en moneda propia (EE.~UU.), la empeora con deuda en divisas.
\end{itemize}
\end{trampa}

\end{columnas}
\end{document}
